\chapter{Internal Models: Forward and Inverse}
\label{ch:internal-models}

\epigraph{%
  To control is to predict.
  To predict is to have a model.
  To have a model is to have learned.
}{}

\section{Forward Models}

A forward model predicts the sensory consequences of a motor command:
\begin{equation}
  \hat{\state}_{k+1} = f(\hat{\state}_k, \control_k),
  \label{eq:ch6:forward-model}
\end{equation}
which is exactly the state propagation used in Volume~I.

\subsection{The MOSAIC Architecture}

The MOSAIC model (Modular Selection and Identification for Control)
\cite{WolpertKawato1998} proposes that the brain maintains \emph{multiple}
paired forward-inverse models, each tuned to a different context:

\begin{lstlisting}[language=Python, caption={MOSAIC-style multiple model architecture.}]
import numpy as np
from typing import List, Tuple

class ForwardInversePair:
    """A paired forward-inverse model for one context."""
    def __init__(self, A: np.ndarray, B: np.ndarray):
        self.A = A  # State transition
        self.B = B  # Input matrix

    def forward_predict(self, x: np.ndarray,
                        u: np.ndarray) -> np.ndarray:
        """Predict next state."""
        return self.A @ x + self.B @ u

    def inverse(self, x: np.ndarray,
                x_desired: np.ndarray) -> np.ndarray:
        """Compute control to reach desired state."""
        return np.linalg.lstsq(self.B,
                               x_desired - self.A @ x,
                               rcond=None)[0]


class MOSAIC:
    """Multiple paired forward-inverse model architecture."""
    def __init__(self, models: List[ForwardInversePair]):
        self.models = models
        self.n_models = len(models)
        self.responsibilities = np.ones(self.n_models) / self.n_models

    def update_responsibilities(self, x: np.ndarray,
                                u: np.ndarray,
                                x_next: np.ndarray) -> None:
        """Update model responsibilities based on prediction errors."""
        errors = np.zeros(self.n_models)
        for i, model in enumerate(self.models):
            prediction = model.forward_predict(x, u)
            errors[i] = np.sum((x_next - prediction)**2)

        # Softmax weighting (lower error = higher responsibility)
        log_resp = -0.5 * errors
        log_resp -= np.max(log_resp)  # numerical stability
        self.responsibilities = np.exp(log_resp)
        self.responsibilities /= np.sum(self.responsibilities)

    def predict(self, x: np.ndarray,
                u: np.ndarray) -> np.ndarray:
        """Weighted prediction across all models."""
        prediction = np.zeros_like(x)
        for i, model in enumerate(self.models):
            prediction += self.responsibilities[i] * \
                         model.forward_predict(x, u)
        return prediction

    def control(self, x: np.ndarray,
                x_desired: np.ndarray) -> np.ndarray:
        """Weighted inverse model output."""
        u = np.zeros(self.models[0].B.shape[1])
        for i, model in enumerate(self.models):
            u += self.responsibilities[i] * \
                 model.inverse(x, x_desired)
        return u
\end{lstlisting}

\section{Inverse Models}

An inverse model computes the motor command required to achieve a desired state:
\begin{equation}
  \control_k = f^{-1}(\state_k, \state_{k+1}^{\text{des}}),
  \label{eq:ch6:inverse-model}
\end{equation}
which corresponds to the feedforward control of Volume~II.

\section{Smith Predictor for Neural Delays}

The sensorimotor system faces delays of 50--200~ms. The Smith predictor
architecture uses a forward model to compensate:
\begin{equation}
  \hat{\state}_{\text{current}} = \hat{\state}(t - \Delta t) +
  \int_{t-\Delta t}^{t} f(\hat{\state}, \control) \, \dd\tau,
  \label{eq:ch6:smith-predictor}
\end{equation}
where $\Delta t$ is the sensory delay. The controller operates on the
\emph{predicted current state} rather than the delayed sensory measurement.

\section*{Exercises}

\begin{enumerate}
  \item \textbf{MOSAIC.} Using the provided code, simulate a system that switches
        between two dynamics (e.g., carrying an empty cup vs.\ a full cup). Show
        how the responsibilities shift.

  \item \textbf{Smith Predictor.} Implement a Smith predictor for a simple pendulum
        with a 100~ms sensory delay. Compare performance with and without prediction.

  \item \textbf{(Research.)} Summarize evidence for multiple internal models in
        human motor control, citing at least three experimental studies.
\end{enumerate}
